\documentclass{article}
\usepackage{graphicx}
\usepackage{spellcheck}

\title{Spellcheck.sty QA Test}
\author{QA Agent}
\date{2026-05-16}

\begin{document}

\maketitle

%% =========================================================
%% TEST 1: \spellerror{word} renders red zigzag underline
%% =========================================================
\section*{Test 1: \textbackslash spellerror renders red zigzag underline}

This word is misspeled and should have a red zigzag underline: \spellerror{misspeled}.

Another example: \spellerror{teh} quick brown fox jumps over the \spellerror{lazzy} dog.

Longer word: \spellerror{accomodation} should also work correctly.

\newpage

%% =========================================================
%% TEST 2: \swarmspellchecktrue / \swarmspellcheckfalse toggle
%% =========================================================
\section*{Test 2: Toggle \textbackslash swarmspellchecktrue / false}

\swarmspellcheckfalse
This \spellerror{misspeled} word should NOT have a zigzag underline (toggle OFF).

\swarmspellchecktrue
This \spellerror{misspeled} word SHOULD have a red zigzag underline (toggle ON).

\newpage

%% =========================================================
%% TEST 3: \spellexport{word} registers correctly
%% =========================================================
\section*{Test 3: \textbackslash spellexport\{word\} registers}

\spellexport{misspeled}
\spellexport{teh}
\spellexport{lazzy}
\spellexport{accomodation}

The words above have been registered via \textbackslash spellexport.  Note:
\textbackslash spellexport only defines a csname — it does NOT automatically
replace in-text occurrences.  This test verifies no errors are raised.

Registered: \spellexport{recieve}
Registered: \spellexport{definately}

\newpage

%% =========================================================
%% TEST 5: Inline helper from spellcheck.py compiles
%% =========================================================
\section*{Test 5: Inline helper (\textbackslash input)}

A misspeled word appears here to trigger the spellchecker.

Another teh example with definately wrong spelling.

\newpage

%% =========================================================
%% TEST 6: \spellchecksummary{N}{total} renders correctly
%% =========================================================
\section*{Test 6: \textbackslash spellchecksummary}

Summary in header area: \spellchecksummary{3}{150}

After disabling: 
\swarmspellcheckfalse
\spellchecksummary{3}{150}

Re-enabled:
\swarmspellchecktrue
\spellchecksummary{5}{200}

\end{document}
